\documentclass[tikz,border=8pt]{standalone}
\usepackage{tikz}
\usepackage{amsmath}
\usepackage{amssymb}
\usepackage{pgfplots}
\pgfplotsset{compat=1.18}
\usepackage{ctex}
\usetikzlibrary{calc,positioning,arrows.meta,matrix}

% LC 62 Unique Paths
% 3x3 网格（行 i=0..2, 列 j=0..2）
% dp[i][j] = dp[i-1][j] + dp[i][j-1]

\begin{document}
\begin{tikzpicture}[
  every node/.style={font=\small},
  cell/.style={draw=gray!70, minimum size=1.2cm, inner sep=0pt, thick, fill=white},
  bordcell/.style={draw=blue!50, minimum size=1.2cm, inner sep=0pt, thick, fill=blue!8},
  rescell/.style={draw=orange!70, minimum size=1.2cm, inner sep=0pt, very thick, fill=orange!8},
  arr/.style={->,>=Stealth,thick,blue!60},
]

%% ── 标题 ──────────────────────────────────────────────────
\node[font=\large\bfseries] at (2.4, 1.4)
  {唯一路径数（LC 62 Unique Paths）};
\node[font=\small, gray] at (2.4, 0.8)
  {$3\times 3$ 网格，dp[i][j] = dp[i\!-\!1][j] + dp[i][j\!-\!1]，从左上到右下};

%% ── 列标签 j=0,1,2（表格上方，y=0.2，不与起点标注冲突）
\foreach \j in {0,1,2} {
  \node[font=\scriptsize, gray] at (\j*1.4+0.7, 0.2) {j=\j};
}
%% ── 行标签 i=0,1,2（表格左侧）
\foreach \i in {0,1,2} {
  \node[font=\scriptsize, gray, anchor=east] at (-0.1, -\i*1.4-0.7) {i=\i};
}
\node[font=\scriptsize, gray, anchor=east] at (-0.1, 0.2) {col→};

%% ── dp 网格（3×3）── cell center: (j*1.4+0.7, -i*1.4-0.7) ──
\foreach \i/\j/\v/\style in {
  0/0/1/bordcell, 0/1/1/bordcell, 0/2/1/bordcell,
  1/0/1/bordcell, 1/1/2/cell,    1/2/3/cell,
  2/0/1/bordcell, 2/1/3/cell,    2/2/6/rescell} {
  \node[\style] (N\i\j) at (\j*1.4+0.7, -\i*1.4-0.7) {\v};
}

%% ── 方向箭头
% 右向
\draw[arr] (N00.east) -- (N01.west);
\draw[arr] (N01.east) -- (N02.west);
\draw[arr] (N10.east) -- (N11.west);
\draw[arr] (N11.east) -- (N12.west);
\draw[arr] (N20.east) -- (N21.west);
\draw[arr] (N21.east) -- (N22.west);
% 下向
\draw[arr] (N00.south) -- (N10.north);
\draw[arr] (N10.south) -- (N20.north);
\draw[arr] (N01.south) -- (N11.north);
\draw[arr] (N11.south) -- (N21.north);
\draw[arr] (N02.south) -- (N12.north);
\draw[arr] (N12.south) -- (N22.north);

%% ── 起点/终点标注（左上角/右下角格子外侧，不占用列标签位置）
\node[font=\tiny\bfseries, blue!70, anchor=east] at (N00.west) {起};
\node[font=\tiny\bfseries, orange!70, anchor=north] at (N22.south) {终点};

%% ── 关键单元格注解：移除右侧带引线的公式框（避免箭头穿过表格中部格子）
%% 单元格转移公式已在下方"填充顺序"说明框中完整列出，此处不重复。

%% ── 分隔线
\draw[gray!30, dashed] (-0.5, -4.5) -- (5.5, -4.5);

%% ── 填充顺序说明（表格下方）
\node[font=\normalsize\bfseries] at (3.5, -5.0) {填充顺序与完整 dp 表};

\node[draw=gray!50, fill=white, rounded corners=3pt, font=\small,
      inner sep=8pt, align=left]
  at (3.5, -6.6)
  {\textbf{初始化：}第0行全为1，第0列全为1（只有一条路可走）\\[4pt]
   \textbf{转移公式：}$\text{dp}[i][j] = \text{dp}[i\!-\!1][j] + \text{dp}[i][j\!-\!1]$\\[4pt]
   \textbf{按行从左到右逐格填入：}\\
   $\text{dp}[1][1]=1+1=2$\quad
   $\text{dp}[1][2]=1+2=3$\quad
   $\text{dp}[2][1]=1+2=3$\quad
   $\text{dp}[2][2]=3+3=\mathbf{6}$\\[4pt]
   \textbf{结果：}机器人从 (0,0) 到 (2,2) 共 $\mathbf{6}$ 条不同路径\\
   \textbf{复杂度：}时间 $O(m \times n)$，空间 $O(m \times n)$ 或 $O(n)$ 滚动数组};

\end{tikzpicture}
\end{document}
